% 324_G6_Z3C_FFGFT_Vergleich_De_ch.tex
% Vergleich: Dougs G(6,Z3C)-Vakuumstruktur und FFGFT
% Numerische Untersuchung der offenen Verbindungsfragen
% Dok. 324, August 2026
%
% Bezugsdokument: D. Matzke, "Vacuum Structure in G(6,Z3C)",
%   ANPA 2026, 13. Juli 2026.
% Verweisdokumente FFGFT:
%   Dok. 049 — Confinement-Formulierung
%   Dok. 160 — alpha_s(m_tau) = 3*xi^{1/4}
%   Dok. 321 — SU(3)_c aus Z3-Trialität [B]
%   Dok. 322 — Hilbertraum-Einbettung
%   Dok. 323 — Weinberg-Winkel aus xi

% =========================================================
% =========================================================
\section*{Zeichenerklärung}
% =========================================================

\subsection*{G(6,$\mathbb{Z}_3\mathbb{C}$)-Symbole nach Matzke (2026)}

\begin{center}
\begin{tabular}{lll}
\toprule
Symbol & Bedeutung & Bemerkung \\
\midrule
$G(6,\mathbb{Z}_3\mathbb{C})$ & Clifford-Algebra, 6 Dimensionen,
  Koeffizientenfeld $\{-1,0,\pm1,\pm i\}$ & $2\equiv-1\pmod{3}$ \\
$e_1,\ldots,e_6$ & Basisvektoren, $e_k^2 = +1$ & \\
$N_k^{\pm}$ & Nilpotente Vektoren, $(N_k^{\pm})^2 = 0$ &
  $N_k^{\pm} = e_a \pm i\,e_b$ \\
$P_{Pk}$ & Idempotente, $P_{Pk}^2 = P_{Pk}$ &
  $P_{Pk} = (\mathbf{1}+i\,e_a e_b)/2$ \\
$V$ & Vakuumoperator, $V = P_{P1}\cdot P_{P2}\cdot P_{P3}$ &
  Rang-1-Projektor \\
$T_k$ & Trine-Operator, $T_k^3 = \mathbf{1}$ &
  $T_k = \mathbf{1}+(\omega-1)P_{Pk}$ \\
$\omega$ & primitiver $\mathbb{Z}_3$-Phasenfaktor &
  $e^{2\pi i/3}$ \\
$I,J,K$ & Kommutative Quaternionen (Grad~4), $I^2=J^2=K^2=+1$ &
  kodieren 3 Raumdimensionen \\
$iE_7$ & Komplexe Volumenform, $(iE_7)^2=+1$ & universeller Phasenoperator \\
\bottomrule
\end{tabular}
\end{center}

\subsection*{FFGFT-Symbole nach Pascher (2023--2026)}

\begin{center}
\begin{tabular}{lll}
\toprule
Symbol & Bedeutung & Wert \\
\midrule
$\xi$ & FFGFT-Geometrieparameter & $4/30000$ \\
$T^4/\mathbb{Z}_3$ & 4D-Torus mit $\mathbb{Z}_3$-Orbifold-Projektion & \\
$\tau$ & Generator der $\mathbb{Z}_3$-Wirkung auf $T^4$ & unitär \\
$P_k$ & $\mathbb{Z}_3$-Projektoren auf $L^2(T^4)$ &
  $\frac{1}{3}\sum_j\omega^{-jk}\tau^j$, $k=0,1,2$ \\
$\mathcal{H}_k$ & Eigenräume: $\mathcal{H}_k = P_k L^2(T^4)$ &
  drei Farbsektoren \\
$\hat{F}_{D_4}$ & $D_4$-Sub-Operator des Spektraloperators &
  $\xi = \lambda_{\min}(\hat{F}_{D_4})$ \\
$C_2(R)$ & Quadratischer Casimir-Operator (Lie-Algebra) &
  $C_2(SU(3)_{\text{fund}}) = 4/3$ \\
$N_{\text{Fourier}}$ & Fourier-Modenzahl des $T^4$-Torus &
  $10^4 = 30000/|\mathbb{Z}_3|$ \\
$N_c$ & Farbfreiheitsgrade, algebraisch notwendig & $3$ \\
$v$ & Higgs-Vakuumerwartungswert & $246\,\text{GeV}$ \\
\bottomrule
\end{tabular}
\end{center}

\subsection*{FFGFT-Dokumentennummerierung}

Dieses Dokument verweist auf den FFGFT-Korpus (Johann Pascher,
2023--2026, ORCID 0009-0000-6518-4064):
\begin{itemize}
  \item \textbf{Dok.~XXX}: Dokument Nr.~XXX des FFGFT-Korpus.
    Vollständige Liste: \texttt{github.com/jpascher/T0-Time-Mass-Duality}
  \item \textbf{RXXX}: Korrektur/Nachtrag Nr.~XXX im Register Dok.~190.
  \item \textbf{[B]}: algebraisch bewiesen \quad
    \textbf{[K]}: aus $\xi$ ableitbar \quad
    \textbf{[S]}: skizziert/plausibel
\end{itemize}
Für diesen Vergleich relevante Dokumente: Dok.~009 (Casimir-Effekt,
Faktor $4/3$), Dok.~049 (Confinement), Dok.~160 ($\alpha_s$),
Dok.~314 ($D_4$-Gitter), Dok.~321 ($SU(3)_c$-Herleitung),
Dok.~322 (Hilbertraum), Dok.~323 (Weinberg-Winkel).


\section*{Zusammenfassung}
% =========================================================

Dieses Dokument vergleicht die Vakuumstruktur von G(6,$\mathbb{Z}_3$C) nach Matzke
mit der FFGFT-Spektraltheorie und untersucht numerisch die offene
Verbindungsfrage: Ist $\xi = 4/30000$ der kleinste Spektralwert des
Vakuumoperators $V$ in G(6)?

Das Ergebnis der numerischen Untersuchung (Python-Skript
\texttt{324\_G6\_FFGFT\_verify.py}):

\begin{itemize}
  \item Das Trine-Theorem ist verifiziert: $T_k^3 = \mathbf{1}$
    bis auf $6{,}5\times10^{-16}$ \textbf{[B]}.
  \item Der Vakuumoperator $V = P_{P1}\cdot P_{P2}\cdot P_{P3}$ ist
    ein Rang-1-Projektor: $V^2 = V$ mit Eigenwerten $\{0^{(7)},\, 1^{(1)}\}$.
  \item $\xi$ ist \emph{kein} direkter Spektralwert von $V$.
    Der Zusammenhang ist \emph{strukturell}, nicht spektral.
  \item Das Trine-Produkt $T_1 T_2 T_3$ ist ein globaler
    $\mathbb{Z}_3$-Operator: $(T_1 T_2 T_3)^3 = \mathbf{1}$,
    Eigenwerte $\{1^{(2)}, \omega^{(3)}, \omega^{2(3)}\}$ \textbf{[B]}.
\end{itemize}

\begin{keyresult}[Statusmarkierung]
\textbf{[K]} aus $\xi$ ableitbar \quad
\textbf{[B]} algebraisch bewiesen \quad
\textbf{[S]} skizziert/plausibel
\end{keyresult}

% =========================================================
\section{Ausgangspunkt: zwei Rahmen, eine Physik}
% =========================================================

\subsection{Matzkes G(6,$\mathbb{Z}_3$C)-Vakuumstruktur}

\textbf{[S]} Matzke konstruiert den Vakuumzustand $V$ der Physik als
algebraisches Objekt in $G(6, \mathbb{Z}_3\mathbb{C})$, einer
Clifford-Algebra über dem Koeffizientenkörper
$\mathbb{Z}_3\mathbb{C} = \{-1, 0, +1, \pm i\}$.
Die zentrale Vakuumrelation lautet
\begin{equation}
  V^2 = -V \quad (\text{in } \mathbb{Z}_3\mathbb{C}\text{-Arithmetik}),
  \label{eq:vakrel}
\end{equation}
wobei $2 \equiv -1 \pmod{3}$. Drei nilpotente Vektorpaare
$N_k^{\pm} = e_a \pm i\,e_b$ mit $(N_k^\pm)^2 = 0$ zerlegen den
$\mathbb{C}^6$-Spinorraum in drei Eigensektoren.
Das Trine-Theorem besagt: für jedes Nilpotent $N$ mit $N^2 = 0$ ist
\begin{equation}
  T = \mathbf{1} + N + (1-\mathbf{1})N \quad \Rightarrow \quad T^3 = \mathbf{1}.
  \label{eq:trine}
\end{equation}

\subsection{Die FFGFT-Spektraltheorie}

\textbf{[K]} In der FFGFT ist $\xi = 4/30000$ der kleinste Eigenwert
des $D_4$-Sub-Operators $\hat{F}_{D_4}$ (Dok.~322):
\begin{equation}
  \xi = \lambda_{\min}(\hat{F}_{D_4}) = \frac{4}{30000}.
  \label{eq:xi_eigen}
\end{equation}
Die $\mathbb{Z}_3$-Projektion $P_k = \frac{1}{3}\sum_{j=0}^2 \omega^{-jk}\tau^j$
auf $L^2(T^4)$ zerlegt den Hilbertraum in drei Eigensektoren
$\mathcal{H}_0, \mathcal{H}_1, \mathcal{H}_2$. Deren bilineare Übergänge
erzeugen die $\mathfrak{su}(3)$-Algebra (Dok.~321, \textbf{[B]}).

% =========================================================
\section{Numerische Verifikation des Trine-Theorems}
% =========================================================

\subsection{Konstruktion der Vakuumstruktur}

Wir verwenden die 8-dimensionale Spinor-Darstellung von G(6) über $\mathbb{C}$,
gegeben durch $\Gamma$-Matrizen als Tensorprodukte der Pauli-Matrizen:
\begin{align}
  \Gamma_1 &= \sigma_x\otimes\mathbf{1}\otimes\mathbf{1}, \quad
  \Gamma_2 = \sigma_y\otimes\mathbf{1}\otimes\mathbf{1}, \quad
  \Gamma_3 = \sigma_z\otimes\sigma_x\otimes\mathbf{1}, \nonumber\\
  \Gamma_4 &= \sigma_z\otimes\sigma_y\otimes\mathbf{1}, \quad
  \Gamma_5 = \sigma_z\otimes\sigma_z\otimes\sigma_x, \quad
  \Gamma_6 = \sigma_z\otimes\sigma_z\otimes\sigma_y.
\end{align}
Die Clifford-Relation $\{\Gamma_i, \Gamma_j\} = 2\delta_{ij}\mathbf{1}_8$
ist erfüllt (numerisch bis $< 10^{-14}$).

Nilpotente über drei Witt-Paare $(0,1)$, $(2,3)$, $(4,5)$:
\begin{equation}
  N_k^+ = \Gamma_{2k} + i\,\Gamma_{2k+1}, \quad
  N_k^- = \Gamma_{2k} - i\,\Gamma_{2k+1}, \quad
  (N_k^\pm)^2 = 0 \quad \textbf{[B]}.
\end{equation}

Idempotente:
\begin{equation}
  P_{Pk} = \frac{\mathbf{1} + i\,\Gamma_{2k}\Gamma_{2k+1}}{2},
  \qquad P_{Pk}^2 = P_{Pk} \quad \text{(numerisch bis } 0\text{)} \quad \textbf{[B]}.
\end{equation}

Vakuumoperator:
\begin{equation}
  V = P_{P1}\cdot P_{P2}\cdot P_{P3}.
  \label{eq:V_def}
\end{equation}

\subsection{Trine-Operatoren und ihr Spektrum}

Mit $\omega = e^{2\pi i/3}$:
\begin{equation}
  T_k = \mathbf{1} + (\omega - 1)\,P_{Pk}.
  \label{eq:trine_def}
\end{equation}

\textbf{[B]} Numerisches Ergebnis:
\begin{equation}
  T_k^3 = \mathbf{1} \quad \text{bis auf } 6{,}5\times10^{-16}.
\end{equation}
Das Trine-Theorem ist in der 8-dimensionalen Spinordarstellung
von G(6) exakt verifiziert.

Das Trine-Produkt:
\begin{equation}
  (T_1 T_2 T_3)^3 = \mathbf{1} \quad \text{bis auf } 1{,}9\times10^{-15}.
\end{equation}
Eigenwerte von $T_1 T_2 T_3$: $\{1^{(2)},\, \omega^{(3)},\, \omega^{2(3)}\}$ --
ein globaler $\mathbb{Z}_3$-Phasenoperator auf dem 8-dimensionalen Spinorraum.

% =========================================================
\section{Spektrum des Vakuumoperators $V$}
% =========================================================

\subsection{Numerische Eigenwerte}

\textbf{[B]} Der Vakuumoperator $V = P_{P1}\cdot P_{P2}\cdot P_{P3}$
hat das Spektrum:
\begin{equation}
  \sigma(V) = \{0^{(7)},\, 1^{(1)}\}.
\end{equation}
$V$ ist ein Rang-1-Projektor. Der einzige nichttriviale Eigenwert ist $\lambda = 1$.

In der $\mathbb{Z}_3\mathbb{C}$-Arithmetik gilt $V^2 = V$ mit
$V^2 \equiv -V \pmod{3}$ (da $V \equiv -V \pmod{3}$ für $V \not\equiv 0$),
was Matzkes Relation $V^2 = -V$ reproduziert.

\subsection{Die Hauptfrage: $\xi$ als Spektralwert?}

\textbf{[K]} $\xi = 4/30000 = 1{,}333\times10^{-4}$ ist \emph{kein}
direkter Spektralwert von $V$. Das Spektrum von $V$ enthält nur $\{0, 1\}$.

Die Verbindung zwischen $\xi$ und $V$ ist \emph{struktureller} Natur:

\begin{table}[h]
\centering
\caption{Struktureller Vergleich G(6,$\mathbb{Z}_3$C) ↔ FFGFT}
\label{tab:vergleich}
\begin{tabular}{lll}
\toprule
Größe & G(6,$\mathbb{Z}_3$C) & FFGFT $T^4/\mathbb{Z}_3$ \\
\midrule
$\mathbb{Z}_3$-Symmetrie & Witt-Zerlegung von $\mathbb{C}^6$ & Orbifold-Projektion auf $T^4$ \\
Drei Sektoren & drei Nilpotent-Paare $N_k^\pm$ & $\mathcal{H}_0, \mathcal{H}_1, \mathcal{H}_2$ \\
$N_c = 3$ & Anzahl Witt-Paare & Ordnung von $\mathbb{Z}_3$ \\
Confinement & $T_R = 0$ (Trialität) & $T_R = 0$, algebraisch [B] \\
Vakuum & Rang-1-Projektor $V$ & Spektralwert $\xi = \lambda_{\min}(\hat{F}_{D_4})$ \\
$SU(3)_c$ & bilineare Übergänge [S] & bewiesen [B] (Dok.~321) \\
Fermionmassen & Nilpotente als Richtungen & Eigenwerte $m_i = r_i\xi^{p_i}v$ [K] \\
Zähler 4 in $\xi$ & 4 Witt-Paare bei $G(8)$ & 4 Torus-Dimensionen \\
Nenner 30000 & $|\mathbb{Z}_3|\times10^4$ (Moden) & $|\mathbb{Z}_3|\times10^4$ (Moden) \\
\bottomrule
\end{tabular}
\end{table}

\subsection{Interpretation: zwei Ebenen derselben Struktur}

\textbf{[S]} Der Vakuumoperator $V$ und der Geometrieparameter $\xi$
beschreiben verschiedene Ebenen derselben geometrischen Struktur:

\begin{itemize}
  \item $V$: algebraische Projektion auf den fermionischen Vakuumstrahl
    in G(6). Spektrum $\{0,1\}$ -- dimensionslos und skaleninvariant.
  \item $\xi$: dimensionslose Skalenverhältnis der Sub-Planck-Geometrie.
    Enthält explizit die Modenzahl $N = 30000$ des $T^4/\mathbb{Z}_3$-Orbifolds.
\end{itemize}

Die Verbindungsformel, die beide Ebenen verknüpft, lautet strukturell:
\begin{equation}
  \xi = \frac{\dim(T^4)}{N_{\text{Moden}}} = \frac{4}{N_c \times 10^4},
  \label{eq:xi_struct}
\end{equation}
wobei $N_c = 3$ (Ordnung von $\mathbb{Z}_3$, algebraisch notwendig, Dok.~321)
und $N_{\text{Moden}} = 10^4$ die Fourier-Modenzahl des $T^4$-Torus ist.
Diese Formel ist keine unabhängige Herleitung, sondern eine Strukturaussage:
Der Zähler~4 zählt die Raumdimensionen (= Anzahl der Witt-Paare in G(8));
der Nenner~30000 zählt die Orbifold-Moden.

% =========================================================
\section{Übereinstimmungen und Divergenzen}
% =========================================================

\subsection{Verifizierte algebraische Übereinstimmungen}

\textbf{[B]} Folgende Strukturen sind in beiden Rahmen identisch und
numerisch verifiziert:

\begin{enumerate}
  \item \textbf{Trine-Theorem:} $T_k^3 = \mathbf{1}$ aus den Nilpotenten
    $N_k^\pm$ -- identisch mit der $\mathbb{Z}_3$-Phasenstukteur der
    FFGFT-Projektoren $P_k = \frac{1}{3}\sum\omega^{-jk}\tau^j$.
  \item \textbf{Drei Sektoren aus einer Symmetrie:}
    G(6) Witt-Zerlegung $\leftrightarrow$ FFGFT $\mathbb{Z}_3$-Eigenraumzerlegung.
    In beiden Rahmen gibt es genau drei Sektoren; ihre bilinearen Übergänge
    erzeugen $\mathfrak{su}(3)$.
  \item \textbf{Confinement als $T_R = 0$:} Matzkes Komplement-Bedingung
    $V(V+1) = 0$ entspricht der Trialitäts-Selektion in Dok.~321.
  \item \textbf{$N_c = 3$ algebraisch erzwungen:}
    Weder G(6) noch FFGFT setzt $N_c$ als freien Parameter.
\end{enumerate}

\subsection{Divergenzen und Ergänzungen}

\begin{table}[h]
\centering
\caption{Offene Brücken zwischen G(6,$\mathbb{Z}_3$C) und FFGFT}
\label{tab:offen}
\begin{tabular}{llll}
\toprule
Punkt & G(6,$\mathbb{Z}_3$C) & FFGFT & Charakter \\
\midrule
Elektroschwache Gruppe & $SU(2)_L$ offen [S] & $SU(2)_L$, $U(1)_Y$ [B] & G(6)-intern \\
Weinberg-Winkel & nicht berechnet & $0{,}2308$ ($-0{,}19\,\%$) [K] & G(6)-intern \\
Leptonmassen & nicht hergeleitet & $m_e, m_\mu, m_\tau < 1\,\%$ [K] & G(6)-intern \\
Neutrinomassen & nicht hergeleitet & $m_{\nu_i}$ aus Fixpunkten [K] & G(6)-intern \\
$\xi$-Verbindung & Casimir-Quotient [K] & $\xi = C_2/N_F$ [K] & \textbf{geschlossen} \\
\bottomrule
\end{tabular}
\end{table}

% =========================================================
\section{Die tatsächliche Verbindungsfrage}
% =========================================================

Die Frage ``ist $\xi$ der kleinste Spektralwert von $V$?
ist durch die numerische Untersuchung und R84 vollständig beantwortet:

\textbf{[K]} $\xi$ ist kein Spektralwert irgendeines G(6)-Operators
-- das ist ein abgeschlossener negativer Befund, keine offene Brücke.
Die Verbindung zwischen G(6) und FFGFT liegt auf einer anderen Ebene:
\begin{equation}
  \xi = \frac{C_2(SU(3)_{\text{fund}})}{N_{\text{Fourier}}}
      = \frac{4/3}{10^4} = \frac{4}{30000}
\end{equation}
verbindet die \emph{algebraische} Struktur von G(6) (Casimir-Invariante
$C_2 = 4/3$, aus $N_c=3$ [B]) mit der \emph{geometrischen} Skala der
FFGFT (Torus-Modenzahl $N_F = 10^4$). Das ist keine G(6)-interne Aussage,
sondern die Brückenformel zwischen beiden Rahmen -- R84.

\textbf{Hinweis zur Elektroschwachen Gruppe:}
$SU(2)_L\times U(1)_Y$ ist in \emph{FFGFT} algebraisch hergeleitet
(Dok.~321 [B]). Die Frage, ob Matzkes $G(6,\mathbb{Z}_3\mathbb{C})$
dieselbe Herleitung intern reproduziert, ist ein \emph{G(6)-internes} Problem
-- keine offene FFGFT-Brücke.

% =========================================================

% =========================================================
\section{$\xi$ aus dem Casimir-Operator von $SU(3)$}
% =========================================================

\subsection{Der quadratische Casimir-Operator}

\textbf{[B]} Der quadratische Casimir-Operator einer Lie-Algebra $\mathfrak{g}$
ist das Gruppeninvariant
\begin{equation}
  C_2(R) = \sum_{a=1}^{\dim\mathfrak{g}} T_a^{(R)} T_a^{(R)},
  \label{eq:casimir_def}
\end{equation}
wobei $T_a^{(R)}$ die Generatoren in der Darstellung $R$ sind.
Für $SU(N_c)$ in der \emph{fundamentalen} Darstellung gilt
\begin{equation}
  C_2(SU(N_c)_{\text{fund}}) = \frac{N_c^2 - 1}{2N_c}.
  \label{eq:casimir_SUN}
\end{equation}
Mit $N_c = 3$ (algebraisch notwendig, Dok.~321, \textbf{[B]}):
\begin{equation}
  C_2(SU(3)_{\text{fund}})
  = \frac{9 - 1}{6} = \frac{8}{6} = \frac{4}{3}.
  \label{eq:casimir_SU3}
\end{equation}

\subsection{$\xi$ als Casimir-Wert pro Fourier-Mode}

\textbf{[K]} Die Fourier-Modenzahl des $T^4$-Torus ist
\begin{equation}
  N_{\text{Fourier}} = \frac{N_{\text{Moden}}}{|\mathbb{Z}_3|}
  = \frac{30000}{3} = 10^4.
  \label{eq:NFourier}
\end{equation}
Damit ergibt sich der FFGFT-Geometrieparameter:
\begin{equation}
  \boxed{\xi = \frac{C_2(SU(3)_{\text{fund}})}{N_{\text{Fourier}}}
  = \frac{4/3}{10^4}
  = \frac{4}{3 \times 10^4}
  = \frac{4}{30000}}
  \label{eq:xi_casimir}
\end{equation}

\textbf{[B]} Diese Formel ist exakt, keine Näherung.
Jede Größe hat eine geometrische/algebraische Begründung:
\begin{itemize}
  \item $C_2(SU(3)_{\text{fund}}) = 4/3$: algebraisch aus $N_c = 3$,
    das wiederum algebraisch notwendig aus der $\mathbb{Z}_3$-Struktur
    (Dok.~321).
  \item $|\mathbb{Z}_3| = 3$: Ordnung der Orbifold-Gruppe (Dok.~321).
  \item $N_{\text{Fourier}} = 10^4$: Fourier-Modenzahl des $T^4$-Torus,
    aus der 4-dimensionalen Topologie und der Auflösungsskala.
\end{itemize}

\subsection{Verbindung zu Dok. 009: Zwei Wege zum Faktor 4/3}

\textbf{[K]} Dok.~009 zeigt, dass der Casimir-Effekt (Vakuumenergie
zwischen Platten) den Faktor $4/3$ aus der QFT liefert:
\begin{equation}
  E_{\text{Casimir}} = -\frac{\pi^2\hbar c}{720\,a^3}\times\frac{4}{3}.
  \label{eq:casimir_eff}
\end{equation}
Dies ist dieselbe Zahl wie $C_2(SU(3)_{\text{fund}}) = 4/3$ — kein Zufall.
Der physikalische Casimir-Effekt (Dok.~009) ist die
Feldtheorie-Manifestation des algebraischen Casimir-Werts der
$SU(3)$-Algebra: In beiden Fällen zählt $4/3$ die
Freiheitsgrade-Gewichtung eines $SU(3)$-invarianten Systems.

\begin{table}[h]
\centering
\caption{Zwei Wege zum Faktor 4/3 in $\xi$}
\label{tab:casimir_wege}
\begin{tabular}{lll}
\toprule
Weg & Quelle & Faktor \\
\midrule
Casimir-Effekt (physikalisch) & QFT, Dok.~009 & $4/3$ aus Vakuumfluktuationen \\
Casimir-Operator (algebraisch) & $SU(3)$, Dok.~321 & $C_2 = (N_c^2-1)/(2N_c)$ \\
\midrule
Ergebnis & beide & $4/3 \to \xi = 4/30000$ \\
\bottomrule
\end{tabular}
\end{table}

\begin{keyresult}[Casimir-Herleitung von $\xi$ \textbf{[K]}]
$\xi = C_2(SU(3)_{\text{fund}})/N_{\text{Fourier}}$ leitet den
Zähler 4 in $\xi = 4/30000$ algebraisch aus dem quadratischen
Casimir-Operator ab. $N_c = 3$ ist algebraisch notwendig (Dok.~321 [B]);
der Faktor $4/3$ erscheint sowohl im algebraischen Casimir als auch
im physikalischen Casimir-Effekt (Dok.~009) — zwei unabhängige
Wege zur selben Zahl.
\end{keyresult}


\section{Fazit und Ausblick}
% =========================================================

\begin{keyresult}[Gesamtstatus Dok.~324 -- $G(6,\mathbb{Z}_3\mathbb{C})$ $\leftrightarrow$ FFGFT]
\textbf{Verifiziert [B]:}
\begin{itemize}
  \item Trine-Theorem: $T_k^3 = \mathbf{1}$ in 8-dim Spinordarstellung
  \item $V$ ist Rang-1-Projektor: $\sigma(V) = \{0^{(7)}, 1^{(1)}\}$
  \item Trine-Produkt $(T_1T_2T_3)^3 = \mathbf{1}$ mit $\mathbb{Z}_3$-Spektrum
  \item Drei Sektoren, bilineare Übergänge, Confinement als $T_R = 0$
  \item $C_2(SU(3)_{\text{fund}}) = 4/3$ aus $N_c = 3$ algebraisch
\end{itemize}
\textbf{Geschlossen [K]:}
\begin{itemize}
  \item $\xi$ ist kein Spektralwert von $V$ (numerisch ausgeschlossen) [K]
  \item $\xi = C_2(SU(3)_{\text{fund}})/N_{\text{Fourier}} = (4/3)/10^4 = 4/30000$
    algebraisch hergeleitet -- R84 [K]
  \item Der Laplace-Operator auf dem Vakuumstrahl ergibt $\Delta = 6\cdot\mathbf{1}$
    (trivial, da $\Gamma_a^2 = \mathbf{1}$) -- kein Zusammenhang mit $\xi$
    als Eigenwert; die Frage ist damit als falsch gestellt abgeschlossen [K]
\end{itemize}
\textbf{Offen [S]:}
\begin{itemize}
  \item $SU(2)_L\times U(1)_Y$ in $G(6,\mathbb{Z}_3\mathbb{C})$
    noch nicht hergeleitet -- 	extbf{G(6)-internes Problem, kein FFGFT-Problem}.
    FFGFT hat $SU(2)_L$ und $U(1)_Y$ bereits aus der Torus-Geometrie
    abgeleitet (Dok.~321 [B]); die Frage ist nur, ob Matzkes
    $G(6,\mathbb{Z}_3\mathbb{C})$ dieselbe Herleitung intern reproduziert
\end{itemize}
\textbf{Bemerkung zum Laplace-Punkt:} Die ursprüngliche Frage
``Gibt es einen natürlichen G(6)-Operator mit $\lambda_{\min} = \xi$?''
ist durch R84 beantwortet: $\xi$ ergibt sich nicht aus einem Spektrum
in G(6), sondern aus dem Casimir-Quotienten $C_2/N_{\text{Fourier}}$ --
einer Verbindung zwischen algebraischer Struktur (G(6)) und geometrischer
Skala (FFGFT-Torus). Das sind zwei verschiedene Ebenen; kein G(6)-Operator
allein kann $\xi$ als Eigenwert haben.
\end{keyresult}

% =========================================================
\section*{Literatur und Quellen}
% =========================================================

\begin{thebibliography}{9}

\bibitem{matzke2026}
Matzke, D. (2026).
\textit{How the Vacuum Builds Space: Nilpotents, Quaternions,
and the Planck Voxel}.
ANPA Conference 2026, 13.~Juli 2026.
Webseite: \texttt{http://www.QuantumDoug.com}

\bibitem{furey2015}
Furey, C. (2015).
Unified theory of ideals.
\textit{Physical Review D} 86, 025018.

\bibitem{furey2018}
Furey, C. (2018).
Three generations, two unbroken gauge symmetries, and
one eight-dimensional algebra.
\textit{Physics Letters B} 785, 84--89.

\bibitem{pascher_korpus}
Pascher, J. (2023--2026).
\textit{FFGFT-Korpus} (Dok.~001--324).
ORCID: 0009-0000-6518-4064.\\
\texttt{https://github.com/jpascher/T0-Time-Mass-Duality}

\bibitem{pascher321}
Pascher, J. (2026). \textit{Dok.~321: Algebraische Herleitung
der $SU(3)_c$-Eichstruktur}. FFGFT-Korpus.

\bibitem{pascher323}
Pascher, J. (2026). \textit{Dok.~323: Herleitung des Weinberg-Winkels
bei $M_Z$}. FFGFT-Korpus.

\bibitem{pdg2024}
Particle Data Group (2024). \textit{Review of Particle Physics}.
\textit{Physical Review D} 110, 030001.

\end{thebibliography}

